\section{Conclusion}
\label{sec:conclusion}

The current \LeanCrypto architecture treats exact execution as primary.
Ordinary probability is cost erasure; a cost bound is a proposition over
supported paths; a natural runtime is a monotone observation of an exact
budget; polynomiality is the final annotation-level certificate; and an
independent operational witness guards admission to the PPT interface.  Typed primitive
handlers, programs, oracle composition, dependent machines, and the UC kernel
all follow this hierarchy.

The practical benefit is not merely uniform terminology.  There is one place
to inspect when checking the probability semantics of an algorithm, one exact
source for each primitive cost, and an explicit proof chain from concrete paths
to a polynomial annotated runtime, followed by a visible operational-admission
obligation for the same run.  Where an abstraction loses information---noncommutative
oracle regrouping, environment-independent possible paths, finite fuel, or
context transport---the required hypothesis is visible in the interface.

This organization establishes a coherent foundation for future cryptographic
proofs while keeping its limits auditable.  Extending the library toward
first-order implementation costs and fully instantiated UC protocols can build
on the same principle: exact execution first, erasure and certificates second,
and host-independent admission only at an explicit boundary.
